% TOP_PROBLEM: {"id": 2302015, "title": "Research Problems in Function Theory — Problem 2.15", "category_id": 8, "category": {"id": 8, "name": "analysis", "display_name": "Analysis", "description": "Limits, continuity, calculus, and function theory.", "slug": "analysis", "order_index": 8, "created_at": "2026-07-31T15:26:25.671Z"}, "status": "open", "problem_number": "AMR-022-2015", "set_id": 13}
% FIRST_POSTED: 2026-10-02
% ENTRY_KIND: statement_only
% UnsolvedMath ID 2302015; source catalogue code AMR-022-2015
% !TeX program = lualatex
% Definition and sources only; this file is not an LLM solution attempt.
\documentclass[11pt,a4paper]{article}
\usepackage[margin=25mm]{geometry}
\usepackage{fontspec}
\setmainfont{lmroman10-regular.otf}[BoldFont=lmroman10-bold.otf,
 ItalicFont=lmroman10-italic.otf,BoldItalicFont=lmroman10-bolditalic.otf]
\usepackage{amsmath,amssymb,mathtools}
\usepackage{unicode-math}
\setmathfont{latinmodern-math.otf}
\AtBeginDocument{\let\setminus\smallsetminus}
\usepackage{xurl,hyperref}
\hypersetup{colorlinks=true,urlcolor=blue}
\setcounter{secnumdepth}{0}
\setlength{\parindent}{0pt}
\setlength{\parskip}{5pt}
\setlength{\emergencystretch}{3em}
\begin{document}
\section*{Research Problems in Function Theory — Problem 2.15}\label{2302015}
{\small UnsolvedMath ID 2302015 · AMR-022-2015 · Analysis · AMR Open Problem Lists}
\subsection{Definitions and mathematical statement}
For an entire function $f:\mathbb C\to\mathbb C$, define
$M_f(r)=\max_{|z|=r}|f(z)|$, $r>0$. Let $f_1,f_2$ be entire functions satisfying $M_{f_1}(r)=M_{f_2}(r)$ for every $r>0$.
Blumenthal's conjecture asks whether they necessarily differ only by rotations and reflections in the input and output planes. More explicitly, must there exist real $\theta,\phi$ such that either
\[
f_2(z)=e^{i\theta}f_1(e^{i\phi}z)
\quad\text{or}\quad
f_2(z)=e^{i\theta}\overline{f_1(e^{i\phi}\overline z)}
\qquad(z\in\mathbb C)?
\]
The conjugation in the second expression is simultaneous in the appropriate input and output, so the resulting function is still holomorphic. Translations are not among the equivalences: they need not preserve circles centred at the origin or the absolute value of the output.
The same question can be restricted to polynomials. The source update gives affirmative results for polynomials with at most four nonzero terms and some real-coefficient situations. The general entire-function question and polynomial cases outside those special classes remain the targets. Equality only near zero or only for large radii is a weaker hypothesis and is not substituted here.
\subsection{Short English statement}
Does the maximum modulus on every centred circle determine an entire function up to rotations and simultaneous reflections?
\subsection{Sources}
\begin{itemize}
\item[S1] ULAM.AI and the UnsolvedMath Contributors, supplied problems.json, record 2302015 (AMR-022-2015), AMR Open Problem Lists. Catalogue identity and source transcription; CC BY 4.0.
\url{https://huggingface.co/datasets/ulamai/UnsolvedMath}.
\item[S2] Hayman - Research Problems in Function Theory (2018), Problem 2.15. Original source indexed by AMR-022-2015.
\url{https://arxiv.org/abs/1809.07200}.
\end{itemize}
\end{document}
